% القسم: المناهج
% الفصل: الاستقراء
% المبحث: الاستقراء البنيوي

\documentclass[../../../include/open-logic-section]{subfiles}

\begin{document}

\olfileid{mth}{ind}{sti}

\olsection{الاستقراء البنيوي}

استعملنا الاستقراء حتى الآن لإثبات نتائج عن جميع الأعداد الطبيعية. لكن
يمكن استعمال مبدأ مقابل له مباشرة لإثبات نتائج عن جميع !!{element} في
مجموعة معرّفة استقرائياً. ويسمى هذا في الغالب \emph{الاستقراء البنيوي}،
لأنه يعتمد على بنية الأشياء المعرّفة استقرائياً.

يعطى التعريف الاستقرائي، بوجه عام، بواسطة: (أ)~قائمة من !!{element}
«الأولية» في المجموعة، و(ب)~قائمة من العمليات التي تنتج !!{element}
جديدة في المجموعة من عناصر قديمة فيها. ففي حالة الحدود السليمة، مثلاً،
تكون الأشياء الأولية هي الحروف. وليس لدينا إلا عملية واحدة؛ وهذه العملية
هي:
\begin{align*}
  o(s_1, s_2) &= [s_1 \circ s_2]
\end{align*}
ويمكنك حتى أن تنظر إلى الأعداد الطبيعية~$\Nat$ نفسها بوصفها معطاة
بتعريف استقرائي: فالشيء الأولي هو~$0$، والعملية هي دالة الخلف~$x + 1$.

لكي نثبت شيئاً عن جميع عناصر مجموعة معرّفة استقرائياً، أي لكي نثبت أن كل
!!{element} في المجموعة له خاصية~$P$، يجب أن:
\begin{enumerate}
\item نثبت أن للأشياء الأولية الخاصية~$P$.
\item نثبت، لكل عملية~$o$، أنه إذا كانت لمعاملاتها الخاصية~$P$، كانت
  للنتيجة أيضاً.
\end{enumerate}
فلكي نثبت شيئاً عن جميع الحدود السليمة، مثلاً، نثبت أنه صحيح لجميع
الحروف، وأنه صحيح لـ$[s_1 \circ s_2]$ متى كان صحيحاً لكل من $s_1$
و$s_2$ على حدة.

\begin{prop}
  في كل حد سليم~$t$، يساوي عدد الرموز $[$ عدد الرموز $]$.
\end{prop}

\begin{proof}
نستعمل الاستقراء البنيوي. فالحدود السليمة معرّفة استقرائياً، وحروفها هي
الأشياء الأولية، و$o$ هي عملية بناء حدود سليمة جديدة من حدود قديمة.
\begin{enumerate}
\item تصح القضية لكل حرف، لأن عدد الرموز $[$ في حرف بمفرده هو~$0$،
  وعدد الرموز $]$ فيه هو أيضاً~$0$.
\item افترض أن عدد الرموز $[$ في $s_1$ يساوي عدد الرموز $]$ في
  $s_1$، وأن الشيء نفسه يصح لـ$s_2$. يكون عدد الرموز $[$ في
  $o(s_1, s_2)$، أي في $[s_1 \circ s_2]$، مساوياً مجموع عددها في
  $s_1$ و$s_2$ مضافاً إليه واحد. ويكون عدد الرموز $]$ في
  $o(s_1, s_2)$ مساوياً مجموع عددها في $s_1$ و$s_2$ مضافاً إليه
  واحد. ومن ثم يساوي عدد الرموز $[$ في $o(s_1, s_2)$ عدد الرموز $]$
  فيه.
\end{enumerate}
\end{proof}

\begin{prob}
  أثبت بالاستقراء البنيوي أنه لا يبدأ أي حد سليم بالرمز~$]$.
\end{prob}

لنقدم برهاناً آخر بالاستقراء البنيوي. القطعة الابتدائية الصحيحة من سلسلة
رموز~$t$ هي أي سلسلة~$s$ توافق $t$ رمزاً فرمزاً عند قراءتها من اليسار،
لكن $t$ أطول منها. فمثلاً، $[a \circ {}$ قطعة ابتدائية صحيحة من
$[a \circ b]$؛ أما $[b \circ {}$ فليست كذلك، لأنهما تختلفان في الرمز
الثاني، ولا $[a \circ b]$، لأنهما متساويتان في الطول.

\begin{prop}\ollabel{prop:initial}
  في كل قطعة ابتدائية صحيحة من حد سليم~$t$ عدد من الرموز $[$ أكبر من
  عدد الرموز $]$.
\end{prop}

\begin{proof}
  نستعمل الاستقراء على~$t$:
  \begin{enumerate}
  \item $t$ حرف بمفرده: عندئذ ليست لـ$t$ أي قطعة ابتدائية صحيحة.
  \item $t = [s_1 \circ s_2]$ لبعض حدين سليمين $s_1$ و~$s_2$. إذا
    كان $r$ قطعة ابتدائية صحيحة من $t$، فثمة عدة احتمالات:
    \begin{enumerate}
    \item $r$ هي $[$ فقط: عندئذ في $r$ رمز $[$ واحد أكثر من الرموز
      $]$.
    \item $r$ هي $[r_1$، حيث $r_1$ قطعة ابتدائية صحيحة من~$s_1$:
      بما أن $s_1$ حد سليم، فإن $r_1$، بحسب فرض الاستقراء، تحتوي على
      رموز $[$ أكثر من الرموز $]$؛ وكذلك $[r_1$.
    \item $r$ هي $[s_1$ أو $[s_1 \circ {}$: بحسب النتيجة السابقة،
      يتساوى عدد الرموز $[$ و$]$ في~$s_1$؛ ولذلك يزيد عدد الرموز $[$
      في $[s_1$ أو $[s_1 \circ {}$ بواحد على عدد الرموز~$]$.
    \item $r$ هي $[s_1 \circ r_2$، حيث $r_2$ قطعة ابتدائية صحيحة
      من~$s_2$: بحسب فرض الاستقراء، تحتوي $r_2$ على رموز $[$ أكثر من
      الرموز $]$. وبحسب النتيجة السابقة، يتساوى عدد الرموز $[$ و$]$
      في~$s_1$. لذلك يكون عدد الرموز $[$ في $[s_1 \circ r_2$ أكبر من
      عدد الرموز~$]$.
    \item $r$ هي $[s_1 \circ s_2$: بحسب النتيجة السابقة، يتساوى عدد
      الرموز $[$ و$]$ في $s_1$، ويصدق الشيء نفسه على~$s_2$. ولذلك
      يوجد في $[s_1 \circ s_2$ رمز $[$ واحد أكثر من الرموز~$]$.
    \end{enumerate}
  \end{enumerate}
\end{proof}

\end{document}
